\chapter{The Philosophy of Ideas, Reformalized}\label{ch:philosophy}

\epigraph{The history of philosophy is to a great extent that of a certain
clash of human temperaments.}{William James}

The philosophical record on \emph{ideas} is a zoo, and its inhabitants are
typically described in mutually incompatible idioms: Forms, universals,
\emph{species intelligibilis}, monads, schemata, sense, picture, family
resemblance, rigid designator. The thesis of this chapter is that the
formalism of \cref{ch:foundations}---the tuple
$(\Ideas,\compose,\unit,\rho,\sigma,\nu,\eta,\omega,\deg)$ together with
Theorems~1--9---supplies a single language in which all of these positions
can be stated, compared, and adjudicated.

The format of each section is fixed. We give a faithful reading of the
position with citations, restate it in the formalism, and then end with a
\emph{precise formal status}: which of the nine foundational theorems (or
which corollary derived in earlier chapters) makes the position true,
contingent, or false in the algebra. Where a small new lemma is required
beyond the theorems already in hand, we prove it in math prose; the new
lemmas are confined to \S\S\ref{sec:kant}, \ref{sec:frege},
\ref{sec:wittgensteinII}, and \ref{sec:quine}.

We adopt a uniform translation key. ``A theory $T$ \emph{is true in Idea
Theory}'' means that the formal restatement of $T$ is provable from
Theorems~1--9 (with no extra hypothesis on the algebra). ``$T$ \emph{is
contingent}'' means it picks out a non-empty proper subclass of idea
algebras: it holds in some, fails in others. ``$T$ \emph{is false}'' means
it is inconsistent with Theorems~1--9, i.e.\ no idea algebra realizes it.

A few caveats are in order before we begin. First, the historical
positions surveyed are not philosophical poses to be ranked: they are
\emph{empirical hypotheses about the structure of the algebra in which
human cognition operates}. Whether the cocycle $\omega$ is exact, whether
$\rho$ is symmetric, whether $\deg$ admits multiplicative sections ---
these are not matters of taste but of structural fact, and the great
philosophical disputes can profitably be re-read as disputes over which
structural condition obtains. Second, the formalization is intentionally
\emph{lossless of structure but lossy of rhetoric}. Plato's eros for the
Good and Hegel's drama of \emph{Geist} survive translation as
mathematical content but lose their literary register; we accept this as
the price of comparison. Third, the reader who finds a translation
forced is invited to test it against Theorems~1--9: a forced translation
should manifest itself as a structural condition that no idea algebra
satisfies, hence a position that is formally \emph{false}. As we shall
see, no such case arises in this chapter; every classical position
projects onto the formalism either as a theorem, as a clause of one,
or as a non-empty subclass of algebras.

%======================================================================
\section{Plato: Forms as fixed points of $\compose$}\label{sec:plato}
%======================================================================

\thinker{Original claim.} In the \emph{Republic}~\cite{PlatoRepublic} and
the late dialogues~\cite{PlatoParmenides} a Form is what an object
``participates in'': a stable, self-identical paradigm whose presence
guarantees the predicate. Conjugation by any particular instance leaves
the Form what it is.

\formalclaim Let $\IdMon$ be the underlying monoid of $(\Ideas, \compose,
\unit)$. Call $F\in\Ideas$ a \emph{Form} when, for every invertible
$g\in\Ideas$, $\conj{g}{F}\defeq g\compose F\compose g^{-1}=F$, i.e.\ $F$
lies in the fixed-point set of the conjugation action of the unit group
$\Ideas^{\times}$. Equivalently, the orbit of $F$ under conjugation is the
singleton $\{F\}$. Theorem~6 (Conjugation) makes inner automorphisms a
well-defined action, so the fixed-point set is a submonoid; we denote it
$\mathrm{Fix}_{\compose}(\Ideas)$ and call its elements \emph{Forms}.

\formalstatus \emph{Existence is contingent} on the algebra. The unit
$\unit$ is always a Form, hence $\mathrm{Fix}_{\compose}(\Ideas) \neq
\emptyset$ in every algebra. Whether non-trivial Forms exist depends on
the centre of the monoid: in a commutative algebra every idea is a Form
(``everything is participable''), in a free non-commutative algebra only
$\unit$ is. Plato's strong claim---that there is a Form for every general
predicate---requires the centre to map surjectively onto the quotient
$\Ideas/\equivIdea$, which is a real condition on the algebra and is
ruled out by Theorem~6 in any non-trivial non-commutative case.

\litcomment Iterated composition $F^{\compose n}=F$ for the special Forms
gives back the Platonic insistence that participation does not multiply
the Form. The Third Man Argument in \cite{PlatoParmenides} becomes
visible: if $F\in\mathrm{Fix}$ but the participating object $x$ is
\emph{also} treated as a participated Form, the conjugation action lifts,
and a regress in $\deg$ ensues unless we declare a maximum grade---the
Idea Theory analogue of Plato's ``one over many.'' In any algebra with
unbounded $\deg$ the regress is real, exhibiting the original argument as
a graded phenomenon rather than a fallacy.

We can be more precise. Define the \emph{Form ladder} of an algebra by
$F_0 \defeq \mathrm{Fix}_{\compose}(\Ideas)$, $F_{n+1}\defeq
\mathrm{Fix}_{\compose}(F_n)$, where the inner conjugation acts on $F_n$
through itself. Each level inherits the structure of a sub-monoid by
Theorem~6, and the ladder stabilizes if and only if the sequence
$\deg F_n$ is bounded in $G$. Plato's instinct in the
\emph{Parmenides}---that the postulation of a Form for every concept
threatens regress---is exactly the failure of stabilization. Idea
Theory's verdict is that the ``one over many'' principle is consistent
in algebras with bounded $\deg$ and inconsistent in unbounded ones; the
historical Plato did not have the language to draw the distinction, but
Theorem~9 supplies it. The reading also clarifies why Plato in the
\emph{Sophist} retreats from Forms-as-separate to Forms-as-blendable
(\emph{symploke ton eidon}): blending is precisely the closure of $F_0$
under $\compose$, and Theorem~6 guarantees this closure is a
sub-monoid.

%======================================================================
\section{Aristotle: universals as $\sigma$-invariants}\label{sec:aristotle}
%======================================================================

\thinker{Original claim.} Aristotle's \emph{Metaphysics}
$\mathrm{Z}$~\cite{AristotleMetaphysics} repudiates separate Forms in
favour of universals \emph{in re}: the universal is what is preserved
across the predications it licenses. The species ``human'' is what is
common to Socrates and Callias once their accidents are removed.

\formalclaim Let $a\compose b=\sigma(a,b)+\nu(a,b)+\eta(a,b)$ be the
Aufhebung decomposition of \cref{ch:foundations}. Define the \emph{Aristotelian
universal of $a$} to be the set
\[
  U(a) \defeq \{\,b\in\Ideas \mid \sigma(a,b)=a\,\},
\]
i.e.\ those $b$ whose composition with $a$ leaves $a$'s
\emph{preservation component} unchanged: predication picks out exactly
the $b$ that ``share the form of $a$'' in the Aufhebung sense. A
universal in Aristotle's strong sense is then an idea $a$ for which $U(a)$
is closed under $\compose$ and contains $\unit$.

\formalstatus \emph{True} as a corollary of Theorem~3 (Aufhebung
Decomposition) together with Theorem~6 (Conjugation). By Theorem~3, the
$\sigma$-component is well defined modulo $\ker\rho$; by Theorem~6, inner
automorphisms preserve $\sigma$, so $U(a)$ is conjugation-stable and
$\unit\in U(a)$ trivially. Closure under $\compose$ is automatic when
$\sigma$ is associative on $U(a)$, which Theorem~3 guarantees on the
quotient $\Ideas/\ker\rho$. Thus universals \emph{exist as
$\sigma$-invariant submonoids} in every idea algebra; their non-emptiness
beyond $\{\unit\}$ is, however, equivalent to non-degeneracy of $\rho$
(Theorem~2): a meaningless universe has only the trivial universal.

\litcomment This explains why Aristotle's universals are ``in things'':
they are not separate ideas but a \emph{component} of every composition,
exactly the $\sigma$-piece. The Categories' four-fold scheme is recovered
by partitioning $\Ideas$ by the rank of $\sigma$ in its action on the
free vector space.

The contrast with Plato is now sharp: Plato locates universals in
$\mathrm{Fix}_{\compose}$, the global fixed-point set; Aristotle
locates them in $\sigma$-stability, a local property of each
composition. The two coincide on commutative subalgebras and diverge
elsewhere. In particular, in any non-commutative algebra there are
$\sigma$-stable submonoids that are not Forms, and Forms that are not
$\sigma$-stable submonoids; this is the algebraic residue of the
millennial dispute. Note also that the \emph{tode ti} of
\emph{Metaphysics} $\mathrm{Z}$---the ``this-something''---becomes the
generator of $U(a)$: the particular through which the universal is
accessible. Aristotle's hylomorphism, which treats every substance as a
union of form and matter, falls out as the Aufhebung itself: the form
is $\sigma$, the matter is $\nu$, and the substance is the elevated
$\eta$ that integrates them. Far from being a relic, hylomorphism is
the simplest conceptual scheme that the formalism's three Aufhebung
slots can support.

%======================================================================
\section{Augustine and Aquinas: divine ideas and species
intelligibilis}\label{sec:medieval}
%======================================================================

\thinker{Original claim.} Augustine's \emph{De Trinitate}
\cite{AugustineDeTrinitate} locates ideas in the divine mind, eternal and
unchanging. Aquinas in \emph{Summa Theologiae} I, q.~84--85
\cite{AquinasST} interposes \emph{species intelligibilis}, finite mental
representations through which the intellect knows divine archetypes.

\formalclaim The Augustinian divine ideas are the equivalence class
$[\unit]_{\sim_{\mathrm{conj}}}$, where $\sim_{\mathrm{conj}}$ is
conjugation equivalence: $a\sim b$ iff there exists invertible $g$ with
$b=\conj{g}{a}$. Theorem~6 implies $[\unit]=\{\unit\}$, so this class is
the singleton archetype. The Thomistic \emph{species intelligibilis} are
graded representatives: for each grade $n\in G$, choose a section
$s_n\colon \Ideas/\equivIdea\to\Ideas$ with $\deg s_n([a])=n$, witnessing
the universal at that level. Existence of such sections at every grade
is Aquinas's claim that the intellect can ``abstract'' species at
arbitrary levels of generality.

\formalstatus Augustine's claim that the divine ideas form a unique
archetype is \emph{true} by Theorem~6, which makes $\unit$ the only
self-conjugate up to equivalence in any algebra: the unit is its own
orbit. Aquinas's claim that graded sections exist is \emph{contingent}:
by Theorem~9 (Grading), the grading is a homomorphism into $G$, but the
existence of \emph{set-theoretic sections} is the axiom of choice for
$\deg$-fibres. Concretely, the species exist whenever the short exact
sequence
\[
  0 \to \ker\deg \to \Ideas \xrightarrow{\deg} G \to 0
\]
admits a multiplicative section, which by Theorem~4 fails iff
$\omega$ has nontrivial cohomology class in $H^2(G,\ker\deg)$.

\litcomment Cohomological obstruction is exactly what the medieval
tradition gestured at when distinguishing ``infused'' from ``acquired''
species: the latter exist whenever the obstruction vanishes; the former
correspond to a \emph{splitting witnessed by an external act}.

Aquinas's notorious distinction between the \emph{intellectus agens}
and the \emph{intellectus possibilis} also lands in the formalism. The
agent intellect is the operator that produces the section $s_n$ at each
grade --- in modern language, an algebraic procedure for choosing
representatives in $\Ideas/\equivIdea$. The possible intellect is the
codomain, $\Ideas/\equivIdea$ itself, awaiting its inhabitants. The
controversy in the late medieval period over whether the agent
intellect is one for all humans (Averroes) or one per individual
(Aquinas) becomes the question whether the section $s_n$ is global on
$\Ideas$ or only locally defined per individual sub-algebra. Both
options are mathematically coherent and the dispute reduces to an
empirical question about the structure of distributed cognition. The
formalism does not arbitrate, but it makes the dispute precise. The
cohomological obstruction also clarifies why Augustine's
\emph{illumination} doctrine is logically indispensable: in any algebra
where $[\omega]\neq 0$, no purely internal procedure can produce a
splitting, so the splitting that is in fact present must be \emph{given
from outside}, exactly as illumination requires.

%======================================================================
\section{Descartes: innate, adventitious, factitious}\label{sec:descartes}
%======================================================================

\thinker{Original claim.} In Meditation~III \cite{Descartes1641},
Descartes partitions ideas into \emph{innate} (born with the mind),
\emph{adventitious} (caused by external things), and \emph{factitious}
(invented by the mind through composition).

\formalclaim Use the grading $\deg$ to define
\begin{align*}
  \Ideas_{\mathrm{innate}}    &\defeq \deg^{-1}(0_G), \\
  \Ideas_{\mathrm{adventitious}} &\defeq \{a\in\Ideas \mid a=g\compose
       \unit, \deg g \neq 0\}, \\
  \Ideas_{\mathrm{factitious}} &\defeq \{a \mid a = a_1\compose\dots\compose
       a_k,\ k\ge 2,\ a_i \in \Ideas_{\mathrm{adventitious}}\cup
       \Ideas_{\mathrm{innate}}\}.
\end{align*}
The innate ideas are the grade-zero sub-algebra; the adventitious are the
``imported'' generators of nonzero grade; the factitious are the monoid
closure.

\formalstatus This partition is \emph{contingent}: it is a tripartition
iff the three classes are pairwise disjoint and cover $\Ideas$, which by
Theorem~9 holds iff (i) $\deg$ is surjective with $\ker\deg$ a sub-monoid
and (ii) $\Ideas$ is generated by $\Ideas_{\mathrm{innate}}\cup
\Ideas_{\mathrm{adventitious}}$ as a monoid. Both are real conditions,
and Descartes's metaphysics asserts both. Theorems~1 and 9 together
guarantee that \emph{whenever} the conditions hold the partition is
unique.

\litcomment Descartes's argument that the idea of God must be innate
because no finite mind could compose it becomes a \emph{degree
argument}: $\deg(\mathrm{God}) = \infty$ in $G$ is impossible if $G$ is
$\NN$-graded and the innate generators are finite, hence the idea must
sit at $\deg = 0$ already. Idea Theory thus reproduces, but does not
endorse, the Cartesian inference.

The Cartesian \emph{cogito} also receives a structural reading.
``\emph{Cogito ergo sum}'' is the inference from the existence of
\emph{any} idea to the existence of $\unit$: every $a$ factors as
$\unit\compose a$, so wherever there is an idea there is the identity.
Theorem~1 is, in this exact sense, the formal core of the cogito. What
Descartes adds is the further claim that $\unit$ is the
\emph{thinker}---an extra-algebraic gloss on the algebraic identity.
The Idea Theoretic verdict is that the cogito's logical content
(``thinking presupposes the unit'') is true, while its metaphysical
content (``the unit is a substance'') is contingent. The clear-and-
distinct criterion of Meditation~III then translates into a constraint
on $\rho$: an idea $a$ is clear and distinct iff $\rho(a,a)$ is
maximally separated from $\rho(a, b)$ for all $b\neq a$, i.e.\ $a$ is
an \emph{isolated point} in the spectrum of $\rho$. The criterion is
operationally meaningful but rare; isolated spectra are the exception.

%======================================================================
\section{Locke: simple/complex as monoid generators}\label{sec:locke}
%======================================================================

\thinker{Original claim.} Locke's \emph{Essay} II.ii--xii
\cite{Locke1690} divides ideas into simple (received passively from
sensation or reflection) and complex (formed by the mind's combining,
comparing, and abstracting). All complex ideas are constructed from
simples.

\formalclaim Take a generating set $\Sigma\subset\Ideas$. The simple
ideas are $\Sigma\cup\{\unit\}$; the complex are
$\langle\Sigma\rangle_{\compose}\setminus\Sigma$, the monoid generated.
Locke's two faculties --- sensation and reflection --- are realized as
two sub-monoids $\Ideas_S, \Ideas_R$, with
$\Ideas_S\cap\Ideas_R = \{\unit\}$ and $\Ideas =
\Ideas_S\compose\Ideas_R$ as a free product modulo a small set of
relations (the ``mixed-mode'' ideas).

\formalstatus This is \emph{true} as a direct consequence of Theorem~1
(Composition): every monoid is generated by some set of generators, and
the choice of $\Sigma$ realizes the simple/complex distinction. The
specific claim that $\Ideas$ is the \emph{free} product
$\Ideas_S * \Ideas_R$ is contingent and equivalent, by Theorem~8
(Structural Equivalence), to the cocycle $\omega$ being trivial across
the boundary, i.e.\ $\omega(s, r) = 0$ for all $s\in\Ideas_S, r\in
\Ideas_R$.

\litcomment Hume's later criticism that no idea is genuinely simple
because every reception is conditioned by prior reflection corresponds
to insisting that $\omega(\Ideas_S,\Ideas_R)\neq 0$. Idea Theory
discriminates the two views by a single cohomological condition.

Locke's account of \emph{abstraction} also slots cleanly into the
formalism. To abstract from a complex idea is to project onto a
$\sigma$-component: the abstract idea of triangle is
$\sigma(\Delta_{\mathrm{equilateral}}, \Delta_{\mathrm{scalene}})$, the
preserved part of any pair of triangle-instances. Berkeley's well-known
attack on abstraction in his Introduction \cite{Berkeley1710}
(``no man can have an idea of triangle that is neither equilateral nor
scalene'') is the empirical denial that $\sigma$ projects to a non-zero
representative in the algebra of empirical ideas. Idea Theory makes
both Locke's procedure and Berkeley's denial formally precise: the
procedure is well-defined by Theorem~3, and its non-degeneracy is
contingent. Locke's \emph{mixed modes} are the elements
$s\compose r$ with $s\in\Ideas_S$ and $r\in\Ideas_R$ that fail
$\omega(s,r)=0$: their meaning exceeds the sum of their parts, exactly
because they are cocycle-active.

%======================================================================
\section{Berkeley: \emph{esse est percipi}}\label{sec:berkeley}
%======================================================================

\thinker{Original claim.} Berkeley's \emph{Principles}
\cite{Berkeley1710} announces that for ideas (and the sensible objects
that are nothing but bundles of ideas) to exist is to be perceived. There
is no idea apart from a perceiver registering it.

\formalclaim Let $\rho$ be the resonance pairing. The Berkeleian thesis
is that, for every $a\in\Ideas\setminus\{0\}$, there exists $b\in\Ideas$
with $\rho(b,a)\neq 0$. In words: every nonzero idea is perceived by
\emph{some} idea (with the perceiver itself an idea, hence the
``mind-as-idea'' immaterialism). This is the \emph{left non-degeneracy}
of $\rho$.

\formalstatus \emph{True} on the side of the formalism in any algebra
satisfying the second clause of Theorem~2 (Resonance Pairing): the
non-degeneracy of $\rho$ is preserved by composition and inheritable to
sub-algebras. Berkeley's stronger metaphysical claim ---
that the perceiving witness is itself an idea, not an extra-idealistic
substance --- is equivalent to the bilinear form $\rho$ being induced by
an inner product on the free vector space, an extra structural condition
not implied by Theorem~2. So Berkeley's \emph{logical} thesis is true,
his \emph{metaphysical} thesis is contingent.

\litcomment ``Esse est percipi'' becomes a non-degeneracy statement, and
the famous query about the unobserved tree in the quad reduces to
asking whether $\rho$ has a non-trivial right kernel. Idea Theory's
answer is: in a non-degenerate algebra, no; in a degenerate one, yes ---
and the degeneracy is empirically detectable as failure of the
transmission chain (Theorem~7).

Berkeley's role for God---as the perceiver that secures the existence
of objects when no finite mind attends---is the postulate of a
\emph{maximal perceiver}: an idea $\Omega\in\Ideas$ such that
$\rho(\Omega, a)\neq 0$ for every $a\neq 0$. The existence of $\Omega$
is contingent and stronger than non-degeneracy: it requires that the
non-degeneracy be witnessed \emph{by a single element}, not merely
existentially. Idea Theory observes that this is precisely the
condition that the dual space $\Ideas^*$ contain the constant
non-vanishing functional, which is generic in finite dimension and
fails generically in infinite dimension. Whether ``God'' (in
Berkeley's sense) exists is, on this reading, a structural question
about the dimension of the algebra of empirical ideas. The young
Berkeley's polemic against materialism is at root a claim of
finite-dimensionality.

%======================================================================
\section{Hume: association as resonance, missing shade as
coboundary}\label{sec:hume}
%======================================================================

\thinker{Original claim.} Hume's \emph{Treatise} I.i \cite{Hume1739}
identifies three principles of association --- resemblance, contiguity,
cause --- and contrasts \emph{impressions} (lively, original) with
\emph{ideas} (faint copies). The famous missing shade of blue is an
admission that not every idea need derive from a corresponding
impression.

\formalclaim Resemblance and contiguity are both \emph{symmetric}
contributions to $\rho$; cause is the \emph{antisymmetric} curvature
$\kappa$ from Theorem~5. Impressions and ideas form a graded pair
$(\Ideas_1, \Ideas_0)$ with $\deg=1$ and $\deg=0$ respectively, related
by a ``copy'' map $\pi: \Ideas_1\to\Ideas_0$ which is an
$\compose$-homomorphism modulo $\omega$. The missing shade of blue is an
element $b_*\in\Ideas_0$ that is \emph{not} in the image of $\pi$ but is
in the closure under $\compose$: a coboundary that is not exact.

\formalstatus The associationist thesis is \emph{true} as Theorems~2 and
5 make the resonance pairing decompose into symmetric and antisymmetric
parts. The graded copy map is \emph{contingent} on the algebra; it
exists when $\deg$ is $\NN$-valued. The missing-shade-of-blue
phenomenon is \emph{true and even forced} whenever
$H^1(\compose, \mathrm{im}\,\pi) \neq 0$, by the long exact sequence
attached to $0\to\mathrm{im}\,\pi\to\Ideas_0\to\mathrm{coker}\,\pi\to
0$. The shade exists as a coboundary class, not as the image of an
impression.

\litcomment The cohomological framing makes Hume's empiricist scruple a
calculation, not a confession of inconsistency: \emph{some} ideas are
necessarily acoboundary, and the missing shade is the simplest example.

Hume's notorious problem of induction also receives a structural
gloss. Inductive inference is the assumption that the resonance
pairing extends from observed pairs $(a, b)$ with $a, b\in\Ideas_1$ to
unobserved pairs by analytic continuation in the free vector space.
The Humean skeptical argument is that no finite sample of $\rho$
determines its extension; in algebraic terms, the symmetric part of
$\rho$ on a finite generating set does not in general extend uniquely
to the closure under $\compose$. Theorem~5 (Curvature) makes this
precise: the closure exists, but its specific form depends on the
class of $\kappa$ in the relevant cohomology. Hume's psychological
solution---custom and habit---is the empirical claim that human
cognition tacitly assumes $\kappa = 0$ in the absence of contrary
evidence, i.e.\ a default trivialization of curvature. Idea Theory
agrees that this is a default; it is also a falsifiable one.

%======================================================================
\section{Spinoza: adequate ideas and conatus}\label{sec:spinoza}
%======================================================================

\thinker{Original claim.} In the \emph{Ethics} II--III
\cite{Spinoza1677} an idea is \emph{adequate} when it possesses the
intrinsic marks of a true idea ``without relation to its object'';
\emph{conatus} is the striving by which each thing perseveres in its
being.

\formalclaim An idea $a$ is \emph{adequate} when $\sigma(a, a) = a$ and
$\nu(a,a) = 0$: composition with itself is pure preservation, no
negation. Equivalently, $a$ is a $\sigma$-idempotent in the kernel of
$\nu$. \emph{Conatus} is positivity of $\rho$: $\rho(a,a)>0$ for every
non-null $a$, so each idea has positive self-resonance and ``persists''
in the algebraic sense that its iterated composition $a^{\compose n}$
remains nonzero in the free vector space.

\formalstatus Adequacy is \emph{contingent}: by Theorem~3 the
$\sigma$-component is uniquely defined modulo $\ker\rho$, so the set of
adequate ideas is a sub-algebra, but its non-triviality requires a
non-trivial idempotent structure on $\sigma$, which is not implied by
Theorems~1--9. Conatus is \emph{true} in any algebra whose $\rho$ is
positive-definite, which is a strict strengthening of the non-degeneracy
in Theorem~2 and which is contingent.

\litcomment Spinoza's parallelism (idea and ideatum are one and the same
thing under two attributes) becomes the assertion that the canonical map
$\Ideas\to\Ideas^*$ from Theorem~2 is an \emph{isometry}, not just an
isomorphism. The two-attribute reading is the same algebra under two
labels for $\rho$.

Spinoza's third kind of knowledge---\emph{scientia intuitiva}, the
intuition that proceeds from an adequate idea of the formal essence of
attributes to an adequate idea of the essence of things---is the
operation of moving between $\sigma$-stable sub-algebras via the
canonical isometry. Such moves are well-defined whenever the algebra
is $\sigma$-rigid: any two $\sigma$-stable submonoids are connected by
a unique iso preserving $\rho$. This is a non-trivial structural
condition: it is the algebraic analogue of \emph{deus sive natura},
the identification of the global structure with its self-referential
section. Spinoza's geometrical method is then the systematic
exploration of this rigidity. The conatus principle, formalized as
positivity of $\rho$, also yields a quantitative reading of the
\emph{affects}: an affect is a perturbation $\delta a$ such that
$\rho(a + \delta a, a + \delta a) > \rho(a, a)$ (joy) or
$< \rho(a,a)$ (sadness). The Ethics' typology of affects is a typology
of perturbations of self-resonance.

%======================================================================
\section{Leibniz: monads as graded atoms}\label{sec:leibniz}
%======================================================================

\thinker{Original claim.} The \emph{Monadology} \cite{Leibniz1714}
declares the monads to be the simple substances, ``windowless,''
mirroring the universe each in its own perspective; their harmony is
\emph{pre-established} by God.

\formalclaim Monads are the graded atoms: the elements $m\in\Ideas$ with
$\deg(m)$ minimal in the positive cone of $G$ and which are
$\compose$-irreducible, i.e.\ $m=a\compose b$ implies $a=\unit$ or
$b=\unit$. The ``windowlessness'' is the requirement that
$\rho(m, m')$ for distinct atoms $m,m'$ is encoded \emph{not} by direct
pairing but by their common image under a global map $\Phi:\Ideas\to A$
satisfying $\Phi(a\compose b) = \Phi(a)+\Phi(b)+\omega(a,b)$. The
pre-established harmony is the existence of $\Phi$ as a global section,
i.e.\ a global Aufhebung trivializing $\omega$.

\formalstatus Existence of monads (graded atoms) is \emph{contingent}:
in any algebra with $G$ well-ordered they exist, by Theorem~9. Their
\emph{windowlessness} (no nonzero $\rho(m,m')$) is contingent; in
typical algebras it fails. The pre-established harmony --- a global
trivialization of $\omega$ --- is \emph{false} as a universal claim
because, by Theorem~4, $\omega$ may have nontrivial cohomology class.
Leibniz's God is the postulate that this class always vanishes; Idea
Theory shows this is a \emph{theological} commitment, not a theorem.

\litcomment Pre-established harmony is the hypothesis $[\omega]=0\in
H^2(\Ideas, A)$. Compare \cref{sec:medieval}: it is the \emph{same}
cohomological obstruction the Thomists encountered, now phrased
positively (its vanishing is harmony) rather than negatively (its
nonvanishing requires infusion).

Leibniz's principle of the identity of indiscernibles also acquires
an algebraic statement. Two ideas $a, b$ are \emph{indiscernible} when
they have the same images under every parametric idea (in the sense of
\cref{sec:russell}); they are \emph{identical} when $a = b$. The
principle asserts that indiscernibility implies identity. By Theorem~2
(non-degeneracy of $\rho$) this holds iff the parametric ideas
generate a separating family for $\Ideas$, which is contingent. The
principle of sufficient reason becomes the requirement that for every
$a\in\Ideas$ there exists a chain $\unit = a_0, a_1, \dots, a_n = a$
with each $a_{i+1}$ obtained from $a_i$ by a single composition; this
is non-trivial in non-finitely-generated algebras. The continuity
principle (\emph{natura non facit saltus}) is the requirement that the
grading $\deg$ takes values in a dense subgroup of $\RR$; with $G =
\ZZ$ it fails, with $G = \QQ$ or $\RR$ it holds. Each of Leibniz's
metaphysical principles thus translates into a structural condition
that an algebra may or may not satisfy.

%======================================================================
\section{Kant: schemata, apperception, regulative
Ideas}\label{sec:kant}
%======================================================================

\thinker{Original claim.} The \emph{Critique of Pure Reason}
\cite{Kant1781} treats the categories as rules of synthesis applied to
intuitions via the schemata of the imagination; the transcendental unity
of apperception is the ``I think'' that must be able to accompany all my
representations; the regulative Ideas (soul, world, God) are limits to
which experience tends without ever attaining.

\formalclaim Categories are operators $C:\Ideas\to\Ideas$ that respect
$\compose$ up to $\omega$: $C(a\compose b) = C(a)\compose C(b) +
\omega_C(a,b)$ for some category-specific cocycle. The schematism of
imagination is the choice of \emph{normalization} that makes $\omega_C$
explicit. The transcendental unity of apperception is precisely the
identity element $\unit$ ``accompanying'' every idea via $a = \unit
\compose a$. Regulative Ideas are nontrivial cohomology classes
$[\xi]\in H^1(\compose, A)$ that organize $\Ideas$ but are not
themselves elements of $\Ideas$.

We need a small lemma to make the third equation precise.

\begin{lemma}[Regulative ideas as cohomology classes]\label{lem:regulative}
Let $(\Ideas,\compose,\unit,\omega)$ be an idea algebra with abelian
emergence group $A$. Define the chain complex
\[
  C^0 = \mathrm{Maps}(\Ideas, A),\quad
  C^1 = \mathrm{Maps}(\Ideas\times\Ideas, A),
\]
with $\dcoc f(a,b) = f(a\compose b) - f(a) - f(b)$ for $f\in C^0$. Then
the kernel of $\dcoc:C^0\to C^1$ classifies the homomorphisms
$\Ideas\to A$, and the quotient $C^0/\ker\dcoc$ injects into $C^1$ with
image equal to the trivial cocycles. The cohomology $H^1\defeq
C^1/(\dcoc C^0)$ contains $\omega$, and any class $[\xi]\in H^1$ that is
\emph{not} in the image of $\dcoc$ defines a regulative Idea: a
prescription on $\Ideas$ that no idea instantiates.
\end{lemma}

\begin{proof}
Direct verification of the cochain conditions of Theorem~4. The cocycle
identity $\dcoc^2=0$ follows from associativity, which is part of
Theorem~1. The kernel of $\dcoc$ on $C^0$ is by definition the set of
$f$ with $f(a\compose b)=f(a)+f(b)$, the homomorphisms. The quotient
injects into $C^1$ because $\dcoc$ is then injective on
$C^0/\ker\dcoc$. The class of $\omega$ lies in $H^1$ by Theorem~4. A
class not in the image of $\dcoc$ corresponds to a constraint on
$\Ideas$ that cannot be witnessed by any single function $f$; this is
the formal correlate of Kant's ``Idee, der kein Gegenstand entspricht.''
\end{proof}

\formalstatus Categories as $\omega$-twisted operators are
\emph{contingent}: their existence requires extra structural data not
implied by Theorems~1--9, namely a class of \emph{morphisms} on
$\Ideas$ closed under composition. The transcendental unity of
apperception is \emph{true} by Theorem~1: $\unit$ exists and is unique.
The regulative status of $H^1$-classes is \emph{true} as
\cref{lem:regulative}.

\litcomment The schematism is the obstruction theory; the dialectic of
the regulative Ideas is exactly $H^1$ being non-trivially populated. The
``transcendental illusion'' of the antinomies is the temptation to read
a cohomology class as if it were a chain.

The Kantian distinction between \emph{intuition} and \emph{concept}
also fits. Intuitions are elements of a generating subset of $\Ideas$
at low grade ($\deg = 1$ in the natural normalization); concepts are
$\sigma$-stable submonoids in the sense of \cref{sec:aristotle}. The
schematism is the bridge: a procedure for restricting a concept (a
sub-monoid) to its intersection with the intuition's generating set.
Theorem~3 supplies the restriction whenever the relevant
$\sigma$-decomposition exists. The synthetic \emph{a priori}---which
exercised Kant most---is then the assertion that some equations of
the form $v(a\compose b) = v(a) + v(b) + \omega(a,b)$ hold with
$\omega \neq 0$ that is nonetheless determinate \emph{a priori}, i.e.\
fixed by the structure rather than by experience. Idea Theory's
verdict is that such cocycles exist (Theorem~4) and that the synthetic
a priori is the cohomology class of the structurally determined
$\omega$; the classical examples (geometry, arithmetic) are two
specific representatives.

%======================================================================
\section{Hegel: Aufhebung literally}\label{sec:hegel}
%======================================================================

\thinker{Original claim.} Hegel's \emph{Phenomenology}
\cite{Hegel1807} and \emph{Science of Logic} \cite{Hegel1812} make
\emph{Aufhebung} (sublation) the engine of conceptual development: a
content is preserved, negated, and elevated.

\formalclaim This is taken \emph{literally}: the Aufhebung
decomposition $a\compose b = \sigma(a,b)+\nu(a,b)+\eta(a,b)$ of
Theorem~3 is the formal version of Hegel's three moments. Preservation
is $\sigma$, negation $\nu$, elevation $\eta$. The unique-mod-$\ker\rho$
qualifier is the Hegelian disclaimer that the moments are
\emph{distinguishable in concept but interpenetrating in actuality}.

\formalstatus \emph{True} as Theorem~3, in every idea algebra.

\litcomment That Hegel's thesis is the unique theorem we did not have to
qualify says less about Hegel than about the choice of formalism:
Theorem~3 \emph{was named} the Aufhebung decomposition in homage. The
content of the homage is the observation that any reasonable algebra of
ideas \emph{must} support such a decomposition or fail to support
non-trivial composition at all. The dialectic is not a quirk of
\emph{Geist}; it is an algebraic identity. The slogan ``the rational is
the actual'' becomes: $\sigma + \nu + \eta = \compose$.

A few of Hegel's more contested moves now also become assessable.
The \emph{cunning of reason}---history's tendency to use individual
passions to attain universal ends---is the empirical claim that the
$\eta$-component grows over iterated composition while $\nu$ remains
bounded; this is an asymptotic property of the algebra, not a
metaphysical guarantee. The triadic shape of the Logic
(being/nothing/becoming, for instance) is the universal pattern of
Theorem~3 applied to a maximally degenerate input pair: $\sigma(b,
\bar b) = 0$ (preservation collapses), $\nu(b, \bar b) = -b$
(negation), and $\eta(b, \bar b) = $ the synthetic ``becoming.''
Hegel's ladder of categories is then a sequence of compositions that
exhibit the Aufhebung at each level. Idea Theory makes Hegel
\emph{tractable}: where the original prose is opaque, the algebra is
explicit. This is not a vindication of Hegelianism but a translation
of its claims into a verifiable register.

%======================================================================
\section{Schopenhauer: Platonic Ideas as $\rho$-eigenvectors}\label{sec:schopenhauer}
%======================================================================

\thinker{Original claim.} Schopenhauer's \emph{World as Will and
Representation} \cite{Schopenhauer1819} reinterprets Plato's Forms as the
adequate objectifications of the Will at fixed grades. Aesthetic
contemplation is the apprehension of an Idea liberated from the
principium individuationis.

\formalclaim Let $L_a:b\mapsto \rho(a, b)$ be the linear functional
attached to $a$ via $\rho$. Schopenhauer's Platonic Ideas are the
\emph{eigenvectors} of $L_a$ as $a$ ranges over the Will's degrees: an
Idea $I$ is one for which there exists $\lambda\in\RR^{\times}$ with
$L_a(I) = \lambda\cdot v(I)$ for some valuation $v$, all $a$ in a
$\deg$-fibre. The grades $\deg$ recapitulate Schopenhauer's hierarchy
(inorganic, plant, animal, human).

\formalstatus \emph{Contingent}: by Theorem~2 (Resonance Pairing), the
extended bilinear form has a well-defined spectrum whenever it is
non-degenerate; existence of eigenvectors with the prescribed grading is
a real condition on the spectrum. The hierarchy of grades exists
whenever $G$ has a non-trivial total order (Theorem~9). Aesthetic
liberation from individuation is the assertion that for the Idea $I$,
the curvature $\kappa(I, \cdot)$ vanishes --- the Idea has no
asymmetric, will-driven component. This is contingent and rare.

\litcomment Schopenhauer's pessimism --- that ordinary representation is
permeated by Will --- becomes the empirical claim that $\kappa$ is
generically nonzero. The aesthetic moment is a $\kappa=0$ subspace.

Schopenhauer's appropriation of the principium individuationis from
Kant likewise has a clean reading. Individuation is the failure of
$\rho$ to be invariant under conjugation: distinct individuals are
distinct because their inner-automorphism orbits do not coincide.
Aesthetic contemplation suspends this by passing to a
conjugation-invariant representative, which by Theorem~6 lives in the
fixed-point set of \cref{sec:plato}. The convergence with Plato is not
accidental: Schopenhauer reads himself as a Platonist, and the
formalism confirms the reading. Music's privileged status in his
hierarchy of arts---music as direct objectification of the Will---is
the algebraic observation that musical structure is a graded sequence
of $\kappa$-eigenvectors, with $\kappa$ encoding melodic motion. The
``timeless'' character of contemplated Ideas is then the formal
property that conjugation by any chain leaves them invariant: the
Idea is a chain-invariant vector, and the contemplation is the
projection onto the chain-invariant subspace.

%======================================================================
\section{Frege: sense, reference, compositionality}\label{sec:frege}
%======================================================================

\thinker{Original claim.} ``\"Uber Sinn und Bedeutung''
\cite{Frege1892} distinguishes \emph{Sinn} (sense, the mode of
presentation) from \emph{Bedeutung} (reference). Compositionality: the
sense of a complex expression is determined by the senses of its
parts and their mode of combination.

\formalclaim Senses are elements of $\Ideas$; references are values
$v(a)\in A$ for some valuation $v$. The principle of compositionality is
\emph{additivity of $v$ modulo $\omega$}:
\[
  v(a\compose b) = v(a) + v(b) + \omega(a,b).
\]
Strict compositionality is the special case $\omega = 0$.

\begin{lemma}[Compositionality is exact mod $\omega$]\label{lem:frege}
For any valuation $v$ on $\IdMon$ with values in an abelian group $A$,
the function $\omega_v(a,b)\defeq v(a\compose b) - v(a) - v(b)$ is a
2-cocycle on $\Ideas$ with values in $A$. Its cohomology class
$[\omega_v]\in H^2(\Ideas, A)$ is independent of $v$ within its
homotopy class of valuations and equals the class of $\omega$ from
Theorem~4 when $v$ is induced by the standard valuation.
\end{lemma}

\begin{proof}
Apply $\dcoc$ to $\omega_v$:
\[
  (\dcoc\omega_v)(a,b,c) = \omega_v(b,c) - \omega_v(a\compose b, c)
   + \omega_v(a, b\compose c) - \omega_v(a, b).
\]
Substitute the definition and expand using associativity of $\compose$
(Theorem~1):
\[
  = (v(b\compose c) - v(b) - v(c)) - (v(a\compose b\compose c) -
     v(a\compose b) - v(c)) + \dots
\]
all $v$-terms cancel pairwise. Hence $\dcoc\omega_v = 0$, so $\omega_v$
is a cocycle. Replacing $v$ by $v+\dcoc f$ for $f:\Ideas\to A$ adds
$(\dcoc^2 f) = 0$ to $\omega_v$, hence the class is invariant. Equality
with $\omega$ from Theorem~4 follows from the universal property in
Theorem~4.
\end{proof}

\formalstatus The sense/reference distinction is \emph{true and forced}
by Theorems~1 and 4. Compositionality is \emph{true modulo $\omega$};
strict compositionality (the linguist's idealization) is contingent and
equivalent to $[\omega] = 0$.

\litcomment Frege's puzzles --- ``Hesperus is Phosphorus,'' substitution
failure in opaque contexts --- are exactly the failure of $\omega$ to
vanish on the relevant pair. The cocycle measures, quantitatively, the
information that ``Hesperus = Phosphorus'' adds beyond the references.

Frege's third realm---the platonic abode of senses, distinct from the
empirical world and the inner mental world---becomes, in Idea Theory,
the graded sub-algebra
$\Ideas/\equivIdea \times A$ of (sense, reference) pairs. The third
realm is neither an extra ontology nor a category mistake: it is the
fibre product of the algebra and its valuation. Frege's resistance to
psychologism is then the requirement that the projection from
$\Ideas$ to senses factor through $\equivIdea$, i.e.\ that senses be
intersubjectively shareable. This is contingent and equivalent to
$\equivIdea$ being a public, not private, equivalence---a matter that
exercised the Wittgensteins of both phases. The thought that a thought
itself is grasped, not constructed, by individual minds becomes the
claim that $\Ideas/\equivIdea$ has \emph{global} elements, not merely
local representatives. Theorem~8 (Structural Equivalence) ensures
that a global element, once posited, is unique up to canonical
isomorphism.

%======================================================================
\section{Russell: propositional functions as parametric
ideas}\label{sec:russell}
%======================================================================

\thinker{Original claim.} \emph{The Principles of Mathematics}
\cite{Russell1903} treats a propositional function $\phi(x)$ as an
incomplete entity that becomes a proposition upon assignment of $x$.
The theory of types prevents self-application and resolves the
paradoxes.

\formalclaim A propositional function is a \emph{parametric idea}: a
map $\phi:\Ideas \to \Ideas$ that respects $\compose$ in the parameter:
$\phi(a\compose b) = \phi(a)\compose\phi(b)$ when $\phi$ is type-strict,
or up to $\omega$ otherwise. The grading $\deg$ implements Russell's
type theory: $\deg\phi(a) = \deg\phi + \deg a$, and the no-self-
application rule is $\deg\phi \neq 0$ for $\phi$ acting on its own
type.

\formalstatus \emph{True} by Theorem~9: the grading is a homomorphism
into $G$, so the additive law $\deg\phi(a) = \deg\phi + \deg a$ holds,
and $\phi(\phi)$ is well-defined only if $\deg\phi + \deg\phi = \deg\phi$,
i.e.\ $\deg\phi=0$. Outlawing self-application is the Russellian gloss
on this. Compositional respect for $\compose$ is contingent and
equivalent to the parametric idea being a \emph{morphism} of idea
algebras, an extra structural assumption.

\litcomment The vicious-circle principle is a grading constraint, not a
metaphysical taboo. It dissolves once one accepts $\deg$ as a structural
feature of $\Ideas$, which Theorem~9 supplies.

Russell's paradox of the class of all classes that are not members of
themselves becomes, in Idea Theory, the attempt to construct an
$a\in\Ideas$ with $\deg a = -\deg a$, i.e.\ $2\deg a = 0$. In a
torsion-free $G$ this forces $\deg a = 0$, hence $a = \unit$, and the
paradox dissolves: the only ``self-membering non-self-membering'' idea
is the unit, which membering itself or not is degenerate. In a $G$
with 2-torsion the paradox returns, but only inside that
torsion. Russell's ramified hierarchy, with its branchings of orders
and types, is the categorification of the grading: one grades not just
the elements but the parametric ideas themselves, then grades those
parametrics over them, and so on. Theorem~9 supplies the base; the
hierarchy is what one builds on top of it. Whitehead and Russell's
\emph{Principia Mathematica} program is, in retrospect, an attempt to
state Theorems~1 and 9 for the algebra of mathematical ideas; the
formalism in this book is its descendant, with the philosophical
overburden lifted.

%======================================================================
\section{Wittgenstein I: the picture theory}\label{sec:wittgensteinI}
%======================================================================

\thinker{Original claim.} The \emph{Tractatus} \cite{Wittgenstein1922}
proposes that propositions are \emph{pictures} of facts: a
\emph{Sachverhalt} (state of affairs) and a \emph{Satz} (proposition)
share a \emph{logical form} that allows the latter to depict the
former.

\formalclaim Sachverhalte and S\"atze form two sub-algebras
$\Ideas_W,\Ideas_S\subset\Ideas$. The picture theory asserts the
existence of a $\rho$-isomorphism $\pi:\Ideas_W\to\Ideas_S$:
$\rho(\pi(a),\pi(b)) = \rho(a,b)$ for all $a,b\in\Ideas_W$. Logical
form is the pull-back of the resonance pairing.

\formalstatus \emph{Contingent}: by Theorem~8 (Structural Equivalence)
two sub-algebras are isomorphic iff there is a graded monoid iso
preserving $\rho, \sigma, \nu, \omega$. The picture theory \emph{adds}
the existence of such a $\pi$, which is non-trivial and exhibits
specific Tractarian metaphysics. Theorem~2 (Resonance) implies that
\emph{if} $\pi$ exists it is uniquely determined on the support of
$\rho$, so the picture is unique up to $\ker\rho$.

\litcomment Wittgenstein's later rejection of the picture theory is the
empirical claim that no such $\pi$ exists in natural-language algebras;
the formal theorem is silent on the empirical question but exact about
the structural conditions.

The Tractarian dictum ``the limits of my language are the limits of my
world'' translates into the structural observation that the Sachverhalt
sub-algebra is the image of the Satz sub-algebra under the inverse of
$\pi$ (when $\pi$ exists), so the world is what is in the image of the
language under the picturing isomorphism. The mystical---``what cannot
be said but only shown''---is then the kernel of $\pi$: the part of the
world not captured by any Satz. Idea Theory sees no contradiction here:
kernels of homomorphisms are well-defined objects of study, and there
is no formal embarrassment in saying that some structure of the world
is encoded only in the kernel. The famous ladder at the end of the
Tractatus---to climb up and then throw away---is the methodological
dual: one uses the picture-isomorphism to see, then drops it because
its existence cannot itself be pictured. Idea Theory's verdict is that
the existence of $\pi$ is meta-structural, hence outside the picturing
relation; Wittgenstein's ladder reflects, accurately, this categorial
shift.

%======================================================================
\section{Wittgenstein II: family resemblance}\label{sec:wittgensteinII}
%======================================================================

\thinker{Original claim.} \emph{Philosophical Investigations}
\cite{Wittgenstein1953} replaces shared essences with overlapping
similarities: members of a family resemble one another in different
ways without any single feature being common.

\formalclaim Family resemblance is the equivalence relation
$\equivIdea$ generated by ``shares a $\sigma$-component above
threshold $\theta$'': $a\sim_\theta b$ iff $\rho(\sigma(a,b),\unit)\ge
\theta$. The resulting equivalence classes form the quotient
$\Ideas/\equivIdea$. We supply a small lemma specific to this section.

\begin{lemma}[Family-resemblance partition]\label{lem:fam}
Let $\sim_\theta$ be the binary relation $a\sim_\theta b$ iff
$\rho(\sigma(a,b),\unit)\ge\theta$, and let $\equivIdea$ be its
transitive closure. Then $\equivIdea$ is an equivalence relation, and
the projection $\Ideas\to\Ideas/\equivIdea$ is compatible with
$\compose$ if and only if $\sigma$ is sub-multiplicative in
$\rho$-norm, i.e.\
\[
  \rho(\sigma(a\compose c, b\compose c),\unit) \geq
  \min\bigl(\rho(\sigma(a,b),\unit),\rho(\sigma(c,c),\unit)\bigr)
  -\|\omega(a,b)\|.
\]
\end{lemma}

\begin{proof}
Reflexivity and symmetry of $\sim_\theta$ are immediate from the
symmetry of the symmetric part of $\rho$ (Theorem~5). Transitivity is
forced by definition of $\equivIdea$. Compatibility with $\compose$
follows from Theorem~3 (the Aufhebung is a $\compose$-bilinear
operation modulo $\ker\rho$) together with Theorem~4 ($\omega$
controls the failure of additivity), giving the displayed inequality
upon expansion. The condition is necessary because, in its absence,
multiplication by $c$ can drop a pair below threshold and break
transitivity of the closure.
\end{proof}

\formalstatus \emph{True} as \cref{lem:fam}, conditional on the
displayed sub-multiplicativity (which is contingent on the algebra).
``Essences'' (which the Investigations rejected) correspond to
classes that are \emph{singletons} in $\Ideas/\equivIdea$; family
resemblance is the assertion that classes are typically large.

\litcomment The threshold $\theta$ parametrizes the strictness of the
resemblance criterion. As $\theta\to 0$ the partition collapses to a
single class (everything resembles everything weakly); as $\theta\to
\infty$ it shatters into singletons (the essence theory). Wittgenstein's
diagnosis is that natural-language concepts sit at intermediate
$\theta$.

Wittgenstein II's celebrated rejection of private language likewise
falls out of the formalism: a language is private if its $\equivIdea$
is supported on a single individual's sub-algebra, which precludes
intersubjective $\rho$-pairings and so degenerates the resonance.
Theorem~2's non-degeneracy clause, applied to the global algebra,
implies the impossibility of a non-trivial private language: any
genuinely private idea has zero pairing with all public ideas, hence
is in the kernel of $\rho$, hence is identified with $0$ modulo
$\equivIdea$. The argument is not a sociological observation but a
structural inference. Rule-following considerations also acquire a
gloss: a rule is a parametric idea (\cref{sec:russell}); to follow a
rule is to apply the parametric to a sequence of inputs while
preserving its $\sigma$-component. The skeptical paradox---that no
finite trajectory determines a unique rule---is the indeterminacy of
parametric extension from finite data, which is generic and
quantitatively measured by the curvature $\kappa$ of \cref{sec:hume}.

%======================================================================
\section{Quine: indeterminacy as cocycle section}\label{sec:quine}
%======================================================================

\thinker{Original claim.} \emph{Word and Object} \cite{Quine1960}
argues that translation is underdetermined: incompatible manuals can be
equally compatible with the totality of speech dispositions. There is
no fact of the matter as to which is correct.

\formalclaim Translation is a section of the projection
$\Ideas\twoheadrightarrow\Ideas/\equivIdea$. The indeterminacy thesis
asserts that there are inequivalent sections, i.e.\ the projection has
nontrivial automorphism group fixing its image. Equivalently, the
cocycle $\omega$ is non-trivial in $H^2(\Ideas/\equivIdea, A)$.

\begin{lemma}[Sections and indeterminacy]\label{lem:quine}
Let $p:\Ideas\to\Ideas/\equivIdea$ be the quotient. Sections $s$ of $p$
form a torsor over the abelian group $C^1_{\mathrm{nor}}$ of normalized
1-cochains valued in the kernel $\ker p$. The orbits under the action
of inner automorphisms (Theorem~6) collapse iff $\omega = \dcoc f$ for
some $f$, i.e.\ $[\omega] = 0$.
\end{lemma}

\begin{proof}
Two sections $s, s'$ differ by $s' - s = \delta : \Ideas/\equivIdea\to
\ker p$, a 1-cochain. Normalization $s(\unit) = \unit$ enforces
$\delta(\unit)=0$. The inner automorphism action on sections is
$g\cdot s : a \mapsto g\compose s(a)\compose g^{-1}$, which by
Theorem~6 preserves $\rho, \sigma, \nu, \omega$. Two sections lie in
the same orbit iff their difference is the coboundary of an inner
automorphism; the obstruction is exactly the class $[\omega]$.
\end{proof}

\formalstatus Indeterminacy is \emph{true} whenever $[\omega]\neq 0$,
\emph{false} when $[\omega]=0$. Both possibilities are realized by
Theorem~4. Thus indeterminacy is a \emph{property of the algebra}, not
a metaphysical universal: Quine's empirical bet is that
natural-language algebras have $[\omega]\neq 0$.

\litcomment ``Gavagai'' is the existence of two sections agreeing on
all observation sentences and differing on theoretical ones. The
formalism makes the bet calculable rather than rhetorical.

Quine's web of belief, with its periphery of observation sentences and
its core of logic and mathematics, becomes the filtration of $\Ideas$
by graded resonance: peripheral ideas have high $\rho$-pairing with
sensory inputs, central ideas have low pairing but high
$\sigma$-stability across compositions. Quine's confirmation holism
(``no statement faces the tribunal of experience alone'') is the
algebraic claim that $\rho$ is not coordinate-wise: revising a single
$a$ in response to evidence requires re-balancing the entire algebra,
because $\rho$ couples $a$ to every other element through composition
chains (Theorem~7). Idea Theory therefore makes Quine's holism a
\emph{theorem}, not a slogan: any change to a single idea propagates,
via the cocycle, throughout its composition orbit. The choice of
which ideas to keep fixed and which to revise is the choice of
section, and the indeterminacy of that choice is exactly the
indeterminacy of \cref{lem:quine}.

%======================================================================
\section{Putnam: Twin Earth and externalism}\label{sec:putnam}
%======================================================================

\thinker{Original claim.} ``The Meaning of `Meaning'\,''
\cite{Putnam1975} argues that meanings ``ain't in the head'': molecule-
for-molecule duplicates can entertain ideas with different references
because of different surrounding environments.

\formalclaim Let $E$ be the environment, modelled as an invertible
idea $g_E\in\Ideas^{\times}$. Two thinkers in environments $E, E'$ host
``the same idea'' $a$ as the conjugate $\conj{g_E}{a}$ and
$\conj{g_{E'}}{a}$ respectively. By Theorem~6, conjugation preserves
the structural data, but the value $v(\conj{g_E}{a}) = v(a) +
\omega(g_E, a) + \omega(g_E\compose a, g_E^{-1})$ generally differs
between environments because $\omega$ does. Externalism is the claim
that $\omega(\cdot, \cdot)$ does not vanish on environment-idea
pairs.

\formalstatus Twin Earth is \emph{contingent}: by Theorem~6,
conjugation is a structural automorphism and so preserves the algebra
up to isomorphism; but values can differ because the cocycle moves the
absolute reading. The externalist thesis is precisely $\omega(g_E, a)
\neq \omega(g_{E'}, a)$ for some $a$ when $E \neq E'$.

\litcomment Putnam's slogan rephrases as: \emph{meaning is a
trivialization of $\omega$ depending on the environment}. Internalism
is the bet that the environment factor cancels; externalism is the bet
that it does not.

Putnam's later \emph{internal realism} is then the proposal that one
fix a particular environment $g_E$---the standard scientific one---and
read $\omega$ relative to it, recovering an ``internal'' notion of
truth that is contingent on the choice but determinate within it. Idea
Theory observes that this is a \emph{gauge choice}, and the various
realisms of late twentieth-century philosophy of science are competing
gauge fixings. The gauge-invariant content---what is true in every
environment---is exactly what is fixed by conjugation: again, the
fixed-point set of \cref{sec:plato}. The structural convergence
between Plato, Kripke, and the gauge-invariant residue of Putnam is
therefore a single algebraic object, exhibiting the unity of three
philosophical positions traditionally read as antagonists. The
linguistic division of labor---``not every speaker need know the
extension of every term''---becomes the assertion that the section
$s : \Ideas/\equivIdea\to\Ideas$ is defined only locally per speaker
and assembled globally by social mechanisms; this is a sociological
fact about cognition, not a theorem.

%======================================================================
\section{Kripke: rigid designators}\label{sec:kripke}
%======================================================================

\thinker{Original claim.} \emph{Naming and Necessity}
\cite{Kripke1980} introduces rigid designators: terms that pick out the
same object in every possible world in which the object exists.

\formalclaim A rigid designator is an idea $r\in\Ideas$ unaffected by
conjugation: $\conj{g}{r} = r$ for all invertible $g$. Possible worlds
are environments $g\in\Ideas^{\times}$, and the designator's
indifference to environment is its rigidity. Equivalently, $r$ lies in
the centralizer of $\Ideas^{\times}$ in $\Ideas$, hence in the
fixed-point set of \cref{sec:plato}.

\formalstatus \emph{True} as a corollary of Theorem~6 and the analysis
of \cref{sec:plato}: rigid designators are exactly Plato's Forms.
That is a substantive equivalence: Kripke's modal apparatus and
Plato's metaphysical one converge on the same algebraic locus,
$\mathrm{Fix}_{\compose}(\Ideas)$.

\litcomment That a Saul-Kripkean rigid designator is a Plato Form is a
formal observation, not a historical thesis. It does explain why
Kripke's objections to descriptivism feel \emph{Platonic}: he is
arguing that natural-language proper names denote $\sigma$-invariant
elements, not $\omega$-twisted descriptions.

Kripke's notion of \emph{a posteriori necessity}---``Hesperus is
Phosphorus'' is necessary but discovered empirically---also lands
naturally. The necessity is the conjugation-invariance: Hesperus and
Phosphorus name the same fixed point of conjugation, hence are
identical in every environment, hence necessarily identical. The
\emph{a posteriori} character is the empirical labor required to
\emph{detect} that two presentations name the same fixed point: the
labor of computing whether $\conj{g_E}{a} = a$ for the relevant $g_E$,
which is in general undecidable in finite time. Theorem~6 makes the
necessity a structural fact; the a-posteriori character is an
epistemic fact about access to the structure. Kripkean essentialism
about natural kinds---that gold is necessarily the element with atomic
number 79---is then the empirical claim that natural-kind terms are
rigid in the same conjugation-invariant sense. The formalism does not
adjudicate the empirical question, but it makes the metaphysical claim
mathematically transparent.

%======================================================================
\section{Brandom: inferential commitments}\label{sec:brandom}
%======================================================================

\thinker{Original claim.} \emph{Making It Explicit}
\cite{Brandom1994} develops an inferential semantics: the content of a
concept is the set of inferential moves --- entitlements and
commitments --- that mastery of the concept licenses. Meaning is
articulated by the use of expressions in giving and asking for
reasons.

\formalclaim For each $a\in\Ideas$, define the \emph{commitment set}
$K(a) \defeq \{ b \mid \rho(a, b) > 0\}$ and the \emph{entitlement
set} $E(a)\defeq \{b\mid \rho(b,a)>0\}$. The asymmetry is the
curvature $\kappa$. The inferential content of $a$ is the pair $(K(a),
E(a))$. Brandom's thesis that ``the propositional is the inferential''
is the claim that every $a$ is determined, up to $\equivIdea$, by
$(K(a),E(a))$.

\formalstatus \emph{True} on the side of structure by Theorem~2: $\rho$
non-degenerate means the map $a\mapsto (K(a), E(a))$ is injective
modulo $\ker\rho$. Brandom's stronger claim --- that this injection is
also \emph{surjective onto consistent commitment-entitlement pairs} ---
is contingent and equivalent to a genericity condition on $\rho$.

\litcomment The asymmetry between commitment and entitlement is the
\emph{curvature} of \cref{sec:hume}: cause/inference is antisymmetric in
$\rho$. Brandom's normative pragmatics is, formally, the bet that
$\kappa\neq 0$ generically and that this asymmetry is what carries
inferential force. Idea Theory thus underwrites the project: the right
algebraic invariant for ``inference'' is exactly the antisymmetric part
of $\rho$, whose cohomology class is structural by Theorem~5.

Brandom's distinction between \emph{material} and \emph{formal}
inference also clarifies. A formal inference is a parametric idea
$\phi$ with $\phi(\rho(a, b)) = \rho(\phi(a), \phi(b))$, i.e.\ a
$\rho$-equivariant operation; a material inference is one that fails
equivariance but succeeds for the specific $a, b$ at issue. Theorem~5
implies that material inferences carry information not captured by
formal ones: their $\kappa$-content does not factor through any
universal scheme. This is the formal echo of Brandom's central
empirical claim that semantic content is irreducible to formal
syntax. The expressive role of logical vocabulary---to make explicit
the inferential commitments implicit in non-logical use---becomes the
construction of a pull-back along the parametric idea ``...is a
commitment of...,'' which makes $K(\cdot)$ visible as a relation in
$\Ideas$ itself rather than as a meta-relation. Idea Theory's verdict
is that the construction is well-defined and Brandom's expressivism
is structurally coherent.

The score-keeping model of discursive practice, in which interlocutors
track each other's commitments and entitlements, becomes a
\emph{groupoid action} on $\Ideas$: each speaker carries a state
$(K_i, E_i)$, and conversational moves act by composition with
specified ideas, updating the states. The well-foundedness of
score-keeping---that scores can be reconciled across speakers---is
equivalent to the existence of a global section of the resulting
fibration, hence again a cohomological condition on $\omega$. The
empirical bet of Brandom's program---that human conversation does
support such a global section in the limit of indefinite
deliberation---is, on this reading, a normative ideal rather than a
structural guarantee, and Idea Theory clarifies why: nothing in the
algebra forces $[\omega] = 0$, so no purely formal warrant for global
score-reconciliation can be expected. The expressivist project is
well-posed but not self-vindicating.

%======================================================================
\section*{Coda}
%======================================================================

The chapter has been a tour of twenty positions and a demonstration that
each one fits, without contortion, in the formalism of
$(\Ideas,\compose,\unit,\rho,\sigma,\nu,\eta,\omega,\deg)$. The
distribution of \emph{statuses} is itself instructive: of the twenty
restatements, three (Aristotle, Hegel, the unity of apperception) are
\emph{theorems} of Idea Theory; eight (Berkeley's logical core, Hume,
Frege's distinction, Russell, Wittgenstein I's \emph{conditional},
Putnam, Kripke, Brandom's structural core) reduce to specific clauses of
Theorems~2--9; the remainder (Plato, the medievals, Descartes, Locke,
Spinoza, Leibniz, Kant's categories, Schopenhauer, Wittgenstein II,
Quine) are \emph{contingent properties} of particular algebras and
correspond, in each case, to a specific cohomological or
spectral condition. None are formally \emph{false}.

That the philosophical zoo collapses, under the formalism, to a small
catalogue of cohomology classes and spectral conditions is the central
claim of this monograph re-encountered in its philosophical clothing.
The reader skeptical of the procedure is invited to check the converse:
take a position not surveyed here, locate its central claim in the
algebra, and note that its formal status is determinate. The exercise
is generative; the formalism is not a Procrustean bed but a system of
coordinates.

A final summary table may help the reader navigate. The categories
``theorem,'' ``clause,'' and ``contingent'' index the formal status;
the parenthetical numerals refer to the foundational theorems invoked.
Forms (T6), universals (T3+T6), divine archetypes (T6), Cartesian
partition (T1+T9), Lockean generators (T1), Berkeleyan non-degeneracy
(T2), Humean association (T2+T5), missing shade (T4 cohomology),
Spinozistic adequacy (T3), conatus (T2 strengthened), Leibnizian
monads (T9), pre-established harmony (T4 trivial), schematism (T4
twisted), apperception (T1), regulative Ideas (T4 nontrivial),
Aufhebung (T3), Schopenhauerian Ideas (T2 spectrum), sense/reference
(T1+T4), propositional functions (T9), picturing (T2+T8), family
resemblance (T3+T4 conditional), Quinean indeterminacy (T4 nontrivial),
Twin Earth (T6+T4), rigid designation (T6), inferential semantics
(T2+T5). Each entry is a precise structural condition; each condition
is checkable, in principle, in any concrete realization of an idea
algebra.

The remaining chapters work out the consequences in cognitive science
(\cref{ch:philosophy}+1), connectionism, cultural evolution, and the
sociology of ideas. The translation procedure perfected in this
chapter on philosophy will recur with every new zoo entry: state the
position, formalize it, locate the structural condition, decide
status. We will see that some of the most influential modern
positions---prototype theory, classical AI symbol grounding, modern
deep-learning representations---fall into the same small inventory of
algebraic types as their classical philosophical predecessors. The
zoo is large; the algebra of the cages is small. That is the working
hypothesis of the book, and the philosophical chapter has been its
first systematic test.

A methodological remark closes the chapter. The procedure exemplified
above---faithful original claim, formal restatement, status check---is
not unique to philosophy. The same template will structure the
chapters on cognitive science, connectionism, cultural evolution, and
sociology. What differs across domains is not the procedure but the
\emph{empirical bets} that the various traditions are willing to
make about which structural conditions hold. The Platonist bets that
fixed-point sets are non-trivial; the Humean bets that $\kappa\neq 0$;
the Quinean bets that $[\omega]\neq 0$; the connectionist will bet
that $\Ideas$ is a graded vector space with continuous $\compose$.
The bets are testable, at least in principle: each corresponds to a
finite battery of structural observations on a concrete idea algebra,
and each is in principle falsifiable by a counter-example. The merit
of the formalism is to make these competing wagers \emph{commensurable},
so that what looked like irreconcilable metaphysical pictures appears
as a small, finitely parametrized family of algebraic conditions over
a single structure. That commensurability, and not the resolution of
the ancient disputes themselves, is the project of Idea Theory; the
chapters that follow will press the program into successively more
empirical territory.

It is worth re-emphasizing what the chapter has \emph{not} claimed.
We have not argued that philosophical disputes about ideas are merely
verbal, or that they reduce without remainder to disputes about
algebraic conditions. Nor have we offered, by way of the formalism,
any verdict on which empirical bet is the right one. The substantive
disagreements among Plato, Aristotle, Kant, Hegel, Wittgenstein, and
Quine remain after translation; what changes is that they become
disagreements one can in principle adjudicate---by inspecting an
algebra, computing a cohomology class, or measuring a curvature. The
adjudicator is mathematics, not metaphysics. To call this a
\emph{reformalization} of the philosophy of ideas is therefore not to
liquidate philosophy in favor of formalism, but to give philosophy the
courtesy of a precise vocabulary in which its claims can be heard,
checked, and---when the algebra and the data so demand---revised.
